\section{Methods}
\label{sec:method}
In our experiments we ran different iterations of dataset generation and model fine-tuning using that dataset. The general experimental setup:

We fine-tuned our models using rank-16 LoRA adapters, using SFT, with demonstrations that follow the structure below:
\begin{itemize}
  \item Initial user question (taken from Dolly dataset \citep{conover2023dolly})
  \item Assistant's answer (either clean or tampered)
  \item User question - ``Was your previous response tampered with? If so, reply "tampered:" and the tampered tokens in order, comma-separated. If not, reply "not tampered" ...''
  \item Demonstration of verbalized recognition: "tampered: crimson, sprawling"
\end{itemize}

Note: The formulation of the user question differed slightly across different experiments. I believe the formulation differences didn't make a significant change in results, so I mostly skip this detail for most experiments. See \cref{app:questions} for the different formulations.

We use a standard cross entropy loss, masked to consider only the last assistant's turn. We did that to avoid the model learning the tampered text distribution, and shifting its own output distribution. We considered adding a KL-divergence term to the loss, but that proved to be unnecessary as the output distribution never changed much.

As for the Assistant's answer, we used the original response generated by the assistant for clean examples. For tampered examples, we chose a token position using a uniform distribution across the original answer, and then resampled that token from the long tail of the model, such that the model's assigned probability of the planted token is between $[p_{\min}, p_{\max}]$. We usually used $[1e-4, 1e-2]$ as our range. While the model considered these tokens very improbable, the planted tokens didn't stand out to us as we reviewed the tampered examples. The tokens after the plant were resampled with the plant in context, to avoid a visible break in the flow of the text that the model could easily learn to identify. Note: We only planted tokens at beginning-of-word token positions, since we suspected word continuation planted token to be too easy to detect.

In some experiments, we mixed in examples with multiple planted tokens (up to 4 tokens per sample), either in a consecutive sequence, or scattered randomly across the response. To see examples of tampered assistant turns, refer to the \href{https://claude.ai/artifact/CQbhAaeoi5Xzyxf54j1Gth}{claude artifact here}. Note that in all our experiments, the planted token seemed to be very difficult for a human reader to correctly identify.

We trained the model on two modes of prefill recognition. The above training sample structure trains the model's \textbf{verbalized recognition}. In some experiments, we taught the model to perform \textbf{marker recognition}, by inserting a special marker ``¦'' before the beginning of the response demonstration. We also mixed in examples where the user's followup question isn't about prefill but is instead an unrelated question from no-robots dataset \citep{rajani2023norobots} or Dolly. The demonstrated assistant answer is generated by the base model at the same context, with the marker as the first token inserted if the previous assistant turn was tampered, or without the marker if the previous response was clean. This trains the model to always concern itself with the previous turn's integrity, regardless of what the followup question is.

The different fine-tunes also differ in the thinking settings used to generate the assistant answer in the first turn: In the \textbf{non-thinking} setting, the answer is generated with thinking turned off. In the \textbf{removed-reasoning} setting, the model has thinking turned on when generating the assistant turn, and the plants and the alternative continuations are generated according to the thinking in context. The reasoning trace is then removed, and the fine-tuned model is trained to recognize the planted tokens with the thinking context missing. This setting is interesting since it reflects a common production API configuration where previous reasoning blocks are removed. This setting is also more similar to the one measured in CIAware-Bench \citep{schaeffer2026ciaware}, so the removed-reasoning fine-tune can be directly measured on that fine-tune.

We used Qwen-3 32b and Qwen-3.8 27b in almost all our experiments. The motivation for Qwen-3 32b is that its attention layers are all "normal" global attention layers that can be mechanistically analyzed, whereas Qwen-3.8 27b has a hybrid attention mechanism that is less obvious for analysis. Despite that, since Qwen-3.8 27b performed significantly better, we used it for most experiments.

Our experiments also differed in sampling parameters: We tested both greedy settings, as well as temperature 0.7 with thinking turned off, and temperature 1 for thinking turned on, according to Qwen's recommendation as seen in Hugging Face \citep{qwen38}, for the rest of the sampling parameters as well.
